\documentclass{article}
\usepackage{amsmath,hyperref}
\title{Retaining healthy carrier loads during MRWS recovery}
\date{}
\begin{document}
\maketitle
In the
\href{https://github.com/automaticdai/MRWS_MultiRobot_Warehouse_Scheduling/tree/941338e9ea4c4eaebb87f8374f80fd576497ad01}{pinned MRWS simulator},
a battery fault on an assigned carrier restarts the order from the
inventory already received by the station. For a multi-carrier order
the native rule sends the surviving carriers home to clear their
unreceived loads. This preserves the required dependency order by doing
some usable pickup work again. The test asks when that work can be
retained without treating a manifest or a release request as a
completed transfer.

\paragraph{Fixture and exact class-count bound.}
The old station has physically accepted a prefix of an ordered
class list, with distinct classes in the worked example. Each
still-usable carrier holds an intact, disjoint consecutive block
of the remaining ordered list, with no overlap with the station
prefix. Its only permitted delivery is that block to the original
station. A new grant is issued only for the exact order after old
delivery powers are disabled; faulted carriers cannot deliver
again during this fixture. The dispatcher receives an authentic
and timely observation of carrier availability and held classes.
Missing classes can be newly picked before their turn. Pick
availability, motion, timing, and interlock implementation are
not supplied by the class-count model.

For $m$ missing classes at this cut and known usable block sizes
$w_1,\ldots,w_n$, the exact worst number of new picks if at most
$k$ further carriers are lost before any retained delivery is
$m$ plus the sum of the largest $\min(k,n)$ sizes. The lower bound
chooses that loss set: work on those carriers can no longer
supply its required classes. The upper bound uses permission to
re-pick and then delivers newly picked classes and surviving
blocks in order. This is a bound on class work conditional on
sufficient picking, routing and station capacity. It is not a
deadline bound or a statement about probabilities. Loss noticed
after a launch that cannot be withdrawn is a different problem.
Overlapping grants, repeated classes without assigned counts,
and noncontiguous partial blocks are not covered.

\paragraph{Executed source calibration.}
The fixture runs the original multi-robot reassignment and
residual-order callbacks, as well as native home, inventory,
station and completion methods with narrow driver stubs. No
network is used. The first critical battery fault leaves a
two-class block unavailable. An earlier class is already held by
the station; two other carriers have two and one outstanding
classes. The native discard path sends both to their homes and
reacquires five classes. The alternative delivers the actual
retained blocks, with only two new pickups. Injecting a second
critical battery loss before either surviving block reaches the
station requires three or four new pickups, depending on the
carrier lost; the discard path still needs five. In every arm,
the station receives the same necessary classes in the same
required order, clears its inventory, and produces exactly one
native order-completion event. The fixture's delivery times and
dispatches are inputs to this callback check, not measurements
from a whole scheduler or movement simulation.

The controller is granted local access to the authoritative
carrier inventory for this replay. The native visualizer feed
is not such a grant: its ordered stream can be absent, and
the emitted classes are not unique physical identities. A useful
operational adapter would need a correctly scoped count of
received stock, still-live carrier loads and revoked old powers,
plus a bounded channel and a source for actual station
settlement. It must not get those from a proof-side private
plant snapshot by default.

\paragraph{What the comparison establishes.}
The original source establishes a real reclaimable cost of
discarding healthy work in this fixture. A direct planner that
already freezes disjoint loads and schedules the same residual
class list ties this strategy. The result is a bounded
calibration for I1, not a comparison against the strongest
published recovery assembly, an independent transfer, or a
compound-fault service guarantee under a movement deadline.
\end{document}
